\begin{foldable} % Overview
    We extensively evaluated our method on BBDF benchmarks.
    We conduct a wide range of ablation studies and visualization to reveal the underlying mechanics of our framework.
    We also provide training details, complementary studies, and extended works can also be found in the Appendix.
\end{foldable}

\subsection{Experimental setup} \label{ssec:04-experiments-setup}
\begin{foldable} % Datasets, architectures, settings
    We evaluate our method across eight benchmark datasets at various difficulty: MNIST~\cite{lecun2002gradient}, USPS~\cite{hull1994database}, SVHN~\cite{netzer2011reading}, FMNIST~\cite{xiao2017fashion}, CIFAR10, CIFAR100~\cite{krizhevsky2009learning}, Tiny-ImageNet~\cite{le2015tiny} (abbreviated: TinyImgN), Imagenette~\cite{howard2020imagenette}.
    We consider three network architectures for the teacher and student, namely LeNet5~\cite{lecun2002gradient}, AlexNet~\cite{krizhevsky2012imagenet}, and ResNet~\cite{he2016deep}.
    These datasets and architectures are the de facto evaluation in BBDFKD methods\cite{wang2021zero, zhang2022ideal, yuan2024data}.
    For detailed training setups, we kindly refer readers to Appendix Sec.~\ref{sec:app01-training}.
\end{foldable}


\subsection{Baselines} \label{ssec:04-experiments-baselines}
\begin{foldable} % Baselines
    We compare our method {\methodname} with the baselines:
    \begin{itemize}
        \item Naive-KD: The student is trained on uniform noise $x\sim\mathcal{X_{\mathcal{U}}}$, where $\mathcal{X}_{\mathcal{U}}=\mathcal{U}\left[0,1\right]^{C\times H\times W}$ and the teacher's hard labels $\hat{y}_{T-\text{Hard}}=T\left(x\right)\in\left\{ 1,\ldots,K\right\}$.
        \item BBDFKD methods: We compare ours with three methods: ZSDB3~\cite{wang2021zero}, IDEAL~\cite{zhang2022ideal}, and DFHL-RS~\cite{yuan2024data}.
    \end{itemize}
    For fair comparisons and where applicable, we used the same/smaller teacher-student network architectures and the same/smaller numbers of synthetic images as in~\cite{wang2021zero, zhang2022ideal, yuan2024data}.
    Their accuracies are obtained from the original papers.
    We repeat each experiment \textbf{five} times and report the average and standard error of student accuracy on the hold-out test set.
\end{foldable}


\begin{table}[ht]
    \caption{Distillation Accuracy (\%) on Simple Datasets}
    \centering
    \begin{tabular}{c|l|l|l|l}
        \hline \textbf{Baseline}
            & \makecell{\textbf{MNIST}  \\ LeNet5  \\ $N$ = 20 K}
            & \makecell{\textbf{USPS}   \\ LeNet5  \\ $N$ = 50 K}
            & \makecell{\textbf{SVHN}   \\ AlexNet \\ $N$ = 50 K}
            & \makecell{\textbf{FMNIST} \\ LeNet5  \\ $N$ = 50 K} \\
        \hline Teacher       &                      99.28 &                      95.47 &                      96.16 &                                 90.90 \\
        \hline Naive-KD      &          96.54$_{\pm .60}$ &          42.05$_{\pm .08}$ &          83.35$_{\pm .57}$ &                    55.11$_{\pm 1.23}$ \\
        ZSDB3                &        96.54$^{[4\times]}$ &            \makecell[c]{-} &            \makecell[c]{-} &                 72.31$^{[1.6\times]}$ \\
        IDEAL                &      96.32$^{[1.3\times]}$ &            \makecell[c]{-} &        84.42$^{[2\times]}$ &                   83.92$^{[2\times]}$ \\
        \textbf{\methodname} & \textbf{98.59}$_{\pm .08}$ & \textbf{94.20}$_{\pm .21}$ & \textbf{95.94}$_{\pm .05}$ & \textbf{83.98}$_{\pm \phantom{0}.65}$ \\
        \hline
        \multicolumn{5}{l}{$N$: synthetic dataset size; Subscript $_{\pm \ldots}$ denotes the standard error;} \\
        \multicolumn{5}{l}{Bracketed superscript $^{[\ldots\times]}$ denotes how much more data the baseline} \\
        \multicolumn{5}{l}{requires; Best results are in \textbf{bold}.} \\
        \multicolumn{5}{l}{DFHL-RS \cite{yuan2024data} did not experiment with these benchmarks.} \\
    \end{tabular}
    \label{tab:standard-simple}
\end{table}

\begin{table}[ht]
    \caption{Distillation Accuracy (\%) on Complex Datasets}
    \centering
    \begin{tabular}{c|l|l|l|l}
        \hline \textbf{Baseline}
            & \makecell{\textbf{CIFAR10}    \\ ResNet18 \\ $N$ = 50K}
            & \makecell{\textbf{CIFAR100}   \\ ResNet18 \\ $N$ = 100K}
            & \makecell{\textbf{TinyImgN}   \\ ResNet18 \\ $N$ = 512K}
            & \makecell{\textbf{Imagenette} \\ ResNet18 \\ $N$ = 50K} \\
        \hline Teacher       &                      95.21 &                       78.36 &                       66.08 &                           91.29 \\
        \hline Naive-KD      &          13.99$_{\pm .23}$ & \phantom{0}2.68$_{\pm .13}$ & \phantom{0}0.75$_{\pm .07}$ &  33.21$_{\pm {\phantom{0}}.39}$ \\
        ZSDB3                &      59.46$^{[1.6\times]}$ &             \makecell[c]{-} &             \makecell[c]{-} &                 \makecell[c]{-} \\
        IDEAL                &        68.82$^{[5\times]}$ &        36.96$^{[20\times]}$ &       27.95$^{[3.9\times]}$ &                 \makecell[c]{-} \\
        DFHL-RS              &      77.86$^{[2.6\times]}$ &       51.94$^{[2.6\times]}$ &             \makecell[c]{-} &                 \makecell[c]{-} \\
        \textbf{\methodname} & \textbf{83.41}$_{\pm .46}$ &  \textbf{52.85}$_{\pm .30}$ &  \textbf{28.03}$_{\pm .22}$ &     \textbf{61.84}$_{\pm 0.77}$ \\
        \hline
        \multicolumn{5}{l}{$N$: synthetic dataset size; Subscript $_{\pm \ldots}$ denotes the standard error;} \\
        \multicolumn{5}{l}{Bracketed superscript $^{[\ldots\times]}$ denotes how much more data the baseline} \\
        \multicolumn{5}{l}{requires; Best results are in \textbf{bold}.} \\
    \end{tabular}
    \label{tab:standard-complex}
\end{table}

\subsection{Standard experiments} \label{ssec:04-experiments-standard}
\begin{foldable} % Overview
    We categorize the datasets into simple and complex datasets based on their difficulty.
    Four simple datasets MNIST, USPS, SVHN, FMNIST all have 10 classes of simple objects, at a $32\times32$ or smaller resolution.
    The complex datasets have more complicated objects or more classes: CIFAR10 and CIFAR100 (10 and 100 classes, $32\times32$), TinyImgN (200 classes, $64\times64$), and Imagenette (10 classes, but high-resolution).
    More dataset details can be found in Appendix Sec.~\ref{ssec:app01-training-datasets}.
\end{foldable}

\begin{foldable}
    We report KD results on simple datasets in Tab.~\ref{tab:standard-simple} and on complex datasets in Tab.~\ref{tab:standard-complex}, and make the following remarks:
    \begin{itemize}
        \item All BBDFKD baselines are better than Naive-KD.
            Notably, our method {\methodname} outperforms all BBDFKD baselines, despite using the least synthetic data.
        \item On complex datasets, NaiveKD has a catastrophic performance drop to near random-guessing, due to having a very imbalanced training set from limited $\mathcal{X}_{\mathcal{U}}$. In fact, we observe that sampling from $\mathcal{X}_{\mathcal{U}}$ plateaus around only 5/100 and 6/200 classes on CIFAR100 and TinyImgN. 
        \item The relative performance gaps between {\methodname} and baselines are relative wider on complex datasets, hinting greater difficulty.
        \item While no other baselines have experimented with Imagenette, {\methodname} moderately achieved 62.78\%.
            With this results, we show that our DIP-KD still works well at high-resolution settings.
    \end{itemize}
\end{foldable}


\subsection{Ablation studies} \label{ssec:04-experiments-ablation}
\begin{foldable}
    In this section, we provide deeper insights into our framework by dissecting the components of our framework, varying the hyperparameters, experimenting with cross-architecture and proxy data settings, and analyzing the embeddings of our image priors.
    We also refer enthusiastic readers to visit Appendix Sec.~\ref{sec:app02-experiments} for complementary studies on KD by compression ratio, alternative optimization objectives, and Appendix Sec.~\ref{sec:app03-extension} for an iterative extension of our method.
\end{foldable}


\subsubsection{Components} \label{sssec:04-experiments-ablation-components}
\begin{table}[ht]
    \caption{Distillation Accuracy (\%) by Inclusion of Components}
    \centering
    \begin{tabular}{l|l|l}
        \hline \makecell{\textbf{Method}} & \makecell{\textbf{MNIST}} & \makecell{\textbf{CIFAR10}} \\
        \hline Naive-KD
            & \phantom{$+$ }96.02$_{\pm .60}$ & \phantom{$+$ }13.99$_{\pm .23}$ \\
        $+$ \textbf{Synthesis:}
            & $+$ \phantom{0}2.27                    & $+$ 64.18 \\
        \quad $^{\blacktriangle}$Hierarchical noise: $\mathcal{D} \gets \{x_\text{hn}\}$
            & \phantom{$+$ 0}1.79$^{\blacktriangle}$ & \phantom{$+$} 59.26$^{\blacktriangle}$ \\
        \quad $^{\blacktriangle}$Nonlinear transform: $\mathcal{D} \gets \{x_\text{nt}\}$
            & \phantom{$+$ 0}0.18$^{\blacktriangle}$ & \phantom{$+$ 0}2.38$^{\blacktriangle}$ \\
        \quad $^{\blacktriangle}$Semantic cutmixing: $\mathcal{D} \gets \mathcal{X}_\mathcal{P}$
            & \phantom{$+$ 0}0.30$^{\blacktriangle}$ & \phantom{$+$ 0}2.54$^{\blacktriangle}$ \\
        $+$ \textbf{Contrastion:} Add $\mathcal{L}_\text{CT}$
            & $+$ \phantom{0}0.16                    & $+$ \phantom{0}3.14 \\
        $+$ \textbf{Distillation:} Add $\mathcal{L}_\text{Soft-KD}$
            & $+$ \phantom{0}0.14                    & $+$ \phantom{0}2.10 \\
        \hline
        Complete \textbf{\methodname}
            & \phantom{$+$ }98.59$_{\pm .08}$        & \phantom{$+$ }83.41$_{\pm .46}$ \\
        \hline
        \multicolumn{3}{l}{$\blacktriangle$ denotes individual contribution of each Synthesis sub-components.} \\
        \multicolumn{3}{l}{Subscript $_{\pm \ldots}$ denotes the standard error.} \\
    \end{tabular}
    \label{tab:abla-components}
\end{table}

\begin{foldable}
    We investigate each component of our framework {\methodname} on MNIST and CIFAR10 with 50K samples each, as representative for simple and complex benchmarks.
    Starting from Naive-KD, we sequentially add each component and report the student accuracy in Tab.~\ref{tab:abla-components}.
    Firstly, Synthesis contributes the most to student accuracy ($+$2.27\% on MNIST, $+$64.18\% on CIFAR10), where the hierarchical noise sub-component has the most impact.
    Other components further improve the student, more apparently on CIFAR10 thanks to its relative difficulty, \ie, $+$3.14\% for Contrastion and $+$2.10\% for Distillation.
    This confirms Synthesis and Contrastion can initialize and optimize data with \textit{diversity}, and Distillation can infuse their inter-class knowledge to the student.
\end{foldable}


\subsubsection{Distilling with proxy data} \label{sssec:04-experiments-ablation-naiveapproaches}
\begin{table}[ht] % tab:abla-naive-approaches
    \caption{Distillation Accuracy (\%) with Proxy Data}
    \centering
    \begin{tabular}{c|c|c|c|c}
        \hline \multicolumn{2}{c|}{\textbf{Training data}} & \multicolumn{3}{c}{\textbf{Evaluation}} \\
        \hline \textbf{Inputs} & \textbf{Labels}          & \makecell{\textbf{MNIST} \\ LeNet5} & \makecell{\textbf{USPS} \\ LeNet5} & \makecell{\textbf{SVHN} \\ AlexNet} \\
        \hline \multirow{2}{*}{MNIST}    & ground-truth & \textbf{99.25}$_{\pm  \phantom{0}.06}$ &          91.88$_{\pm           1.56}$ &          20.77$_{\pm \phantom{0}.29}$ \\
                                         & pseudo       & \textbf{99.06}$_{\pm  \phantom{0}.05}$ &          93.94$_{\pm \phantom{0}.30}$ &          41.18$_{\pm           3.27}$ \\
        \hline \multirow{2}{*}{USPS}     & ground-truth &          77.92$_{\pm            2.72}$ & \textbf{95.90}$_{\pm \phantom{0}.48}$ &          19.36$_{\pm \phantom{0}.62}$ \\
                                         & pseudo       &          96.78$_{\pm  \phantom{0}.42}$ & \textbf{94.97}$_{\pm \phantom{0}.15}$ &          57.52$_{\pm           1.81}$ \\
        \hline \multirow{2}{*}{SVHN}     & ground-truth &          64.10$_{\pm            1.29}$ &          62.12$_{\pm           1.68}$ & \textbf{95.97}$_{\pm \phantom{0}.08}$ \\
                                         & pseudo       &          98.59$_{\pm  \phantom{0}.18}$ &          93.95$_{\pm \phantom{0}.37}$ & \textbf{96.48}$_{\pm \phantom{0}.04}$ \\
        \hline $\mathcal{X}_\mathcal{N}$ &              &          88.05$_{\pm  \phantom{0}.76}$ &          27.29$_{\pm \phantom{0}.10}$ &          84.42$_{\pm \phantom{0}.48}$ \\
               $\mathcal{X}_\mathcal{U}$ & pseudo       &          96.54$_{\pm  \phantom{0}.60}$ &          42.05$_{\pm \phantom{0}.08}$ &          83.35$_{\pm \phantom{0}.57}$ \\
               $\mathcal{X}_\mathcal{P}$ &              & \textbf{98.62}$_{\pm  \phantom{0}.17}$ & \textbf{94.14}$_{\pm \phantom{0}.52}$ & \textbf{95.84}$_{\pm \phantom{0}.10}$ \\
        \hline
        \multicolumn{5}{l}{Subscript $_{\pm \ldots}$ denotes the standard error.} \\
        \multicolumn{5}{l}{The top-3 results from each target evaluation (column) are in \textbf{bold}.} \\
    \end{tabular}
    \label{tab:abla-naive-approaches}
\end{table}

\begin{foldable} % Other naive approaches
    Thanks to the digits datasets (MNIST, UPSP, SVHN) having shared classes (digits 0--9), we also explore using out-of-domain datasets as proxy data as another naive approach.
    We evaluate students trained with different combinations of input images (real training sets, noises, and our image priors) and labels (ground truth from real training sets, or pseudo-labels given by the teacher).
    We use 50K images for each synthetic source, and the full training set for each real dataset.
    For fair comparisons, we only include the \textit{Synthesis} component (\ie, $S_{0}$ accuracy) for our method, and reuse the teachers from standard experiments.

\end{foldable}

\begin{foldable} % Discussion on tab:abla-naive-approaches
    We emphasize that using real training sets of shared classes is \textit{unfair} to our settings (which starts without any).
    We summarize our results in Tab.~\ref{tab:abla-naive-approaches} and conclude: 
    \begin{itemize}
        \item Undoubtedly, whenever the teacher's original training set is available (in-domain), the student benefits the most.
        \item Distilling in an adaptation-like manner (\eg, from MNIST to USPS), \ie, querying the teacher on proxy datasets provides a basic improvement compared to using drop-in labels.
            However, even though having access to high-quality images, it \textbf{fails} to give satisfactory performance.
            We hypothesize that this is caused by domain gap \cite{ben2010theory}.
            Hence, using public datasets as a drop-in proxy can potentially harm KD performance.
        \item Rather, instead of \textit{exploitation} on an \textit{unsure} domain, we should priortize \textit{exploration} on a \textit{diverse} domain.
            This is confirmed by our results with image priors $\mathcal{X}_\mathcal{P}$ (bottom row), approximating the in-domain settings with less than 1\% difference.
    \end{itemize}
    These insights suggest the \textbf{universal} diversity of our image priors for kickstarting KD on various benchmarks.
\end{foldable}


\subsubsection{Embeddings analysis --- Generation} \label{sssec:04-experiments-ablation-embeddings-generation}
\begin{foldable}
    To ensure that our image priors are diverse and semantically meaningful, we sample random synthetic images and extract their embeddings from the backbone of the ResNet18 teacher network pre-trained on CIFAR10.
    Its intermediate features have a dimension of 512, where we further embed with UMAP~\cite{mcinnes2018umap} with default parameters to two-dimensional for visualization.
    To quantitatively analyze synthetic images, we evaluate four distribution metrics from generative models~\cite{naeem2020reliable}:
    \begin{itemize}
        \item \textit{Precision}: the fraction of synthetic samples that is in the neighborhood of at least one real sample.
            High precision indicates synthetic data are captured within the support of real data, but may be \textit{biased} by real outlier samples.
        \item \textit{Recall}: the fraction of real samples that is in the neighborhood of at least one synthetic sample.
            High recall indicates synthetic data can capture the support of real data, but may be \textit{biased} by an overestimated synthetic manifold.
        \item \textit{Density}: the average normalized count of synthetic samples in the neighborhood of each real sample.
            High density indicates \textit{fidelity} like precision but is less \textit{biased}, and may be greater than 1 when synthetic data concentrates around real modes.
        \item \textit{Coverage}: the fraction of real samples whose neighbourhoods contain at least one fake sample.
            High coverage indicates \textit{diversity} like recall but is less \textit{biased}.
    \end{itemize}
\end{foldable}

\begin{figure}[ht] % fig:embeddings-universal
    \centering
    \includegraphics[width=\columnwidth]{./figures/04-experiments-vizembed-2d-objects.png}
    \caption{Feature embeddings of complex datasets, noise, and image priors.}
    \label{fig:embeddings-universal}
\end{figure}

\begin{table}[ht]
    \caption{Distribution Metrics (\%) on Complex Datasets}
    \label{tab:abla-embeddings-coverage-objects}
    \centering
    \begin{tabular}{c|c|c|c|c}
        \hline  & CIFAR10 & CIFAR100 & TinyImgN & Imagenette \\
        \hline \textbf{Density} & & & & \\
        $\mathcal{X}_\mathcal{U}$ & \phantom{0}20.00 &           100.65 &           108.65 &           101.20 \\
        $\mathcal{X}_\mathcal{N}$ & \phantom{0}20.00 &           107.98 & \phantom{0}90.87 &           103.64 \\
        $\mathcal{X}_\mathcal{P}$ & \phantom{0}49.14 & \phantom{0}96.20 & \phantom{0}96.00 & \phantom{0}92.14 \\
        \hline \textbf{Coverage} & & & & \\
        $\mathcal{X}_\mathcal{U}$ & \phantom{000}.00 & \phantom{000}.10 & \phantom{000}.06 & \phantom{000}.05 \\
        $\mathcal{X}_\mathcal{N}$ & \phantom{000}.00 & \phantom{000}.10 & \phantom{000}.06 & \phantom{000}.05 \\
        $\mathcal{X}_\mathcal{P}$ & \phantom{00}7.33 & \phantom{0}60.79 & \phantom{0}63.88 & \phantom{0}78.48 \\
        \hline \textbf{Precision} & & & & \\
        $\mathcal{X}_\mathcal{U}$ &           100.00 &           100.10 &           100.00 &           100.00 \\
        $\mathcal{X}_\mathcal{N}$ &           100.00 &           100.10 &           100.00 &           100.00 \\
        $\mathcal{X}_\mathcal{P}$ & \phantom{0}80.83 & \phantom{0}97.76 & \phantom{0}99.29 & \phantom{0}99.10 \\
        \hline \textbf{Recall} & & & & \\
        $\mathcal{X}_\mathcal{U}$ & \phantom{000}.00 & \phantom{000}.09 & \phantom{000}.06 & \phantom{000}.04 \\
        $\mathcal{X}_\mathcal{N}$ & \phantom{000}.00 & \phantom{000}.10 & \phantom{000}.06 & \phantom{000}.05 \\
        $\mathcal{X}_\mathcal{P}$ & \phantom{0}66.92 & \phantom{0}96.75 & \phantom{0}96.86 & \phantom{0}95.07 \\
        \hline
    \end{tabular}
\end{table}

\begin{foldable}
    We explore four complex datasets CIFAR10, CIFAR100, TinyImgN, and Imagenette, which contain rich diversity and semantics.
    We also include uniform noise $\mathcal{X}_{\mathcal{U}}$, Gaussian noise $\mathcal{X}_{\mathcal{N}}$, and our image priors $\mathcal{X}_{\mathcal{P}}$.
    Apparently in Fig.~\ref{fig:embeddings-universal}, $\mathcal{X}_{\mathcal{U}}$ and $\mathcal{X}_{\mathcal{N}}$ form a distant cluster away from the regions of images, suggesting that their semantics are limited and dissimilar to the natural objects.
    In contrast, $\mathcal{X}_{\mathcal{P}}$ widely spreads among the mutual regions of real images, hinting that they share similar semantics.
    The statistics in Tab.~\ref{tab:abla-embeddings-coverage-objects} further confirm our image priors have superior fidelity and diversity.
    On a side note, precision is clearly overestimated, while CIFAR10 metrics drop a bit as this is the backbone's domain knowledge, where it discriminate much harsher.
    Overall, these results explains why image priors $\mathcal{X}_{\mathcal{P}}$ excels in BBDF settings.
\end{foldable}

\subsubsection{Embeddings analysis --- Contrastion} \label{sssec:04-experiments-ablation-embeddings-contrastion}
\begin{foldable}
    To highlight the effect of the Contrastion step on image priors, we visualize the 50K embeddings before-and-after Contrastion on CIFAR10 in Fig.~\ref{fig:embeddings-cifar10}.
    Having fitted, the backbone can proficiently recognize the images and embed them into ten distinctive clusters (colored gray).
    Where image priors reside, the Contrastion step further improves the coverage and expands them to neighbor regions as well.
\end{foldable}

\begin{figure}[ht] % fig:embeddings-universal
    \centering
    \includegraphics[width=\columnwidth]{./figures/04-experiments-vizembed-2d-cifar10-resnet18.png}
    \caption{Feature embeddings of image priors $\mathcal{X}_\mathcal{P}$ on CIFAR10 before and after the Contrastion step, with real data and noise sources.}
    \label{fig:embeddings-cifar10}
\end{figure}

\begin{foldable}
    Upon closer inspection of the randomly-picked raw images (Fig.~\ref{fig:abla-contrastion-images}), we observe that the contrastion optimization has overlay the image priors with faint but non-trivial patterns.
    In most images, the residual can reinforce the existing patterns.
    We also observe a significant increase in coverage and recall, with a small trade-off in density (Tab.~\ref{tab:abla-embeddings-coverage-cifar10}).
    We conclude the Contrastion step empirically improves data diversity.
\end{foldable}

\begin{figure}[ht] % fig:abla-contrastion-images
    \begin{minipage}{\columnwidth}
        \centering
        \includegraphics[width=0.99\columnwidth]{./figures/experiments-abla-contrastion-cifar10-resnet18-before.png}
        {\footnotesize \textbf{(a)} Before Contrastion}
        \vspace{0.5\baselineskip}
    \end{minipage}
    \begin{minipage}{\columnwidth}
        \centering
        \includegraphics[width=0.99\columnwidth]{./figures/experiments-abla-contrastion-cifar10-resnet18-after.png}
        {\footnotesize \textbf{(b)} After Contrastion}
    \end{minipage}
    \caption{Contrastion update to image priors. (Best viewed when zoomed in.)}
    \label{fig:abla-contrastion-images}
\end{figure}

\begin{table}[ht]
    \caption{Distribution Metrics (\%) during Contrastion}
    \label{tab:abla-embeddings-coverage-cifar10}
    \centering
    \begin{tabular}{l|c|c|c|c}
        \hline  & \textbf{Density} & \textbf{Coverage} & \textbf{Precision} & \textbf{Recall} \\
        \hline
        Before Contrastion  & 70.46 & \phantom{0}6.13 & 91.69 & 82.18 \\
        After Contrastion   & 66.93 &           14.66 & 91.68 & 94.00 \\
        \hline
    \end{tabular}
\end{table}


\subsubsection{Increased synthetic dataset size} \label{sssec:04-experiments-ablation-nsamples}
\begin{foldable}
    While it is unsurprising that we achieved better KD performance when using more a larger number of image priors $N$ (Fig.~\ref{fig:abla-nsamples}), more important discussions are: which is the optimal cut-off point, and how it corresponds to the relative difficulty of each dataset.
    We observe that for the digits datasets (MNIST, UPSP, SVHN), the student fits easily even with minimal data ($N$ = 10K), whereas it struggles a bit more on FMNIST and CIFAR10.
    For CIFAR100, TinyImgN, and Imagenette, the student still stably improves even after we have reached $N$ = 100K or more, suggesting that increasing $N$ is a viable option.
\end{foldable}

\begin{figure}[ht] % fig:abla-nsamples
    \centering
    \includegraphics[width=\columnwidth]{./figures/experiments-abla-nsamples.pdf}
    \caption{Distillation accuracy (\%) by number of image priors $N$.}
    \label{fig:abla-nsamples}
\end{figure}


\subsubsection{Cross-architecture KD} \label{sssec:04-experiments-ablation-crossarch}
\begin{foldable}
    An under-investigated setting in BBDFKD methods is distilling from the teacher to a student of different architecture.
    To fully discriminate the knowledge gap, we evaluate three architectures with large capacity difference: LeNet5, AlexNet, and ResNet18; on three benchmarks with intermediate difficulty: FMNIST, SVHN, and CIFAR10.
    We reuse the same teacher networks from the standard experiments for the reference results on the diagonal, and report the results in Tab.~\ref{tab:abla-crossarch}.
    We observe that the student performs best when it has a homogeneous architecture to the teacher.
    Otherwise, accuracy scales with capacity, and a medium-size student can still often give moderate results.
    Evidently, our method can work reasonably well in cross-architecture settings.
\end{foldable}

\begin{table}[ht]
    \caption{Distillation Accuracy (\%) of Cross-Architecture KD}
    \centering
    \begin{tabular}{c|l|l|l}
        \hline
            \makecell{\textbf{Network}}
            & \makecell{\textbf{FMNIST}  \\ LeNet5}
            & \makecell{\textbf{SVHN}    \\ AlexNet}
            & \makecell{\textbf{CIFAR10} \\ ResNet18} \\
        \hline
        Teacher  &                        90.90 &                           96.16 &                        95.21 \\
        \hline
        LeNet5   & \textbf{83.98}$_{\pm \phantom{0}.65}$ &                    64.44$_{\pm 1.83}$ &                    32.01$_{\pm 3.18}$ \\
        AlexNet  &          78.98$_{\pm \phantom{0}.92}$ & \textbf{95.94}$_{\pm \phantom{0}.05}$ &                    61.30$_{\pm 0.69}$ \\
        ResNet18 &                    81.69$_{\pm 1.58}$ &          95.91$_{\pm \phantom{0}.04}$ & \textbf{83.41}$_{\pm \phantom{0}.46}$ \\
        \hline
        \multicolumn{4}{l}{Subscript $_{\pm \ldots}$ denotes the standard error.} \\
        \multicolumn{4}{l}{Best results are in \textbf{bold}.} \\
    \end{tabular}
    \label{tab:abla-crossarch}
\end{table}

